\appendix
\chapter{Practical Implementation and Code Structure}\label{ch:implementation}

\section{Repository Overview}\label{sec:repo-overview}
% High-level description of the project structure:
% - src/: Core algorithms (features, gram, mismatch, evaluation)
% - scripts/: CLI entry points for inference
% - experiments/: Scripts for reproducing LOO and MIL experiments
% - thesis/: LaTeX source for this document

\section{Core Modules}\label{sec:core-modules}

\subsection{\texttt{mismatch.py}}\label{subsec:code-mismatch}
% Details the recursive generation of the mismatch neighbourhood using the
% epigenetic alphabet, heavily leveraging Python's \texttt{itertools} and
% \texttt{@lru\_cache} for memoisation.

\subsection{\texttt{features.py}}\label{subsec:code-features}
% Explains the integration with \texttt{scikit-learn}'s \texttt{CountVectorizer}
% and the parallel chunking logic that returns sparse CSR matrices.
% Emphasises the stateless design achieved via pre-enumerated vocabularies.

\subsection{\texttt{gram.py}}\label{subsec:code-gram}
% Describes the MKL weight generation, dense Gram matrix normalisation via
% NumPy broadcasting, and the block-partitioned multiplication for both
% symmetric and asymmetric matrices.

\subsection{\texttt{evaluation.py} and \texttt{plotting.py}}\label{subsec:code-eval}
% Outlines the OC-SVM fitting via \texttt{scikit-learn}, sequence-to-patient
% score aggregation, and the decoupling of computation from matplotlib
% plotting logic.

\section{Reproducing the Experiments}\label{sec:reproduction}
% Command-line instructions on how to run the pipeline.
% Example: python experiments/run_loo_experiment.py --data-dir path/to/data --tumor-file path/to/tumor.fa
% Mention the role of the seed argument to ensure exact reproducibility.
